\documentclass[10pt,a4paper,sans]{moderncv}
\moderncvstyle{classic}
\moderncvcolor{blue}
\usepackage[utf8]{inputenc}
\usepackage[scale=0.82]{geometry}

% moderncv compose \today en anglais par défaut ; format français sans charger babel.
\renewcommand{\today}{\number\day\space\ifcase\month\or janvier\or février\or mars\or avril\or
mai\or juin\or juillet\or août\or septembre\or octobre\or novembre\or décembre\fi\space\number\year}

\name{Alexandro}{Disla}
\title{Lettre de Motivation}
\address{Rue Saint-Louis Jeanty \#14, Route de Frère}{Port-au-Prince, Haïti}{}
\phone[mobile]{(+509) 4148-3700}
\email{alexandrodisla@hotmail.com}
\extrainfo{GitHub: \url{https://github.com/AD0791}}

\begin{document}

\recipient{Équipe de Recrutement --- Acted Haïti}{Port-au-Prince, Haïti\\Poste : Assistant(e) MEAL --- Réf : ASSISTMEAL \_ 2610}
\date{\today}
\opening{Madame, Monsieur,}
\closing{Je vous prie d'agréer, Madame, Monsieur, l'expression de mes salutations distinguées,}
\enclosure[Ci-joint]{Curriculum Vitae, références, attestation C.T.P.E.A, copie de la carte d'identité}

\makelettertitle

Je vous adresse ma candidature au poste d'Assistant(e) MEAL d'Acted Haïti. J'apporte cinq ans de suivi-évaluation en ONG en Haïti --- la santé et le VIH à la Caris Foundation International, l'eau et l'assainissement chez HANWASH, l'éducation chez Anseye Pou Ayiti --- précédés de huit ans de suivi d'indicateurs de projets publics au Ministère de la Planification. Je réside à Port-au-Prince ; le français et le créole haïtien sont mes langues maternelles.

La collecte et la qualité des données sont le cœur de mon travail. J'ai construit les instruments de collecte sur KoboToolbox et ODK, avec logiques de saut et contraintes de validation, pour que l'erreur soit arrêtée à la saisie plutôt que corrigée trois semaines plus tard, et j'ai formé les agents de collecte à leur usage, guides techniques à l'appui. À la Caris Foundation, j'ai encadré techniquement les agents de collecte d'un programme multi-sites, établi et supervisé les protocoles de DQA en remontant les chiffres rapportés jusqu'aux registres sources, et produit les indicateurs MER destinés à l'USAID via DATIM : je connais le rythme et les contrôles de cohérence qu'impose un rapportage bailleur.

Votre fiche précise qu'en l'absence d'une équipe dédiée à la base de données, l'Assistant(e) MEAL entretient les systèmes d'information du programme, y compris le rapportage sur les indicateurs 3 Zéros. C'est là que je serai le plus utile. J'exerce aussi comme ingénieur de données : j'administre des bases relationnelles, je nettoie et dédoublonne les listes de bénéficiaires, et j'ai automatisé en Python et SQL des rapports d'activité, ce qui a réduit de 40 \% le temps de traitement manuel. Une base propre à chaque projet et des tableaux de bord Power BI tenus à jour chaque mois, c'est un travail que je sais faire, pas seulement spécifier.

Vous verrez que mon expérience dépasse ce que demande le poste ; je le choisis parce que la collecte, la base et l'analyse sont le travail que je veux continuer à faire. Je préfère aussi poser deux réserves franchement. Je n'ai pas encore opéré moi-même un mécanisme de retours d'information : j'apprendrai les procédures du MRA d'Acted plutôt que de prétendre les connaître. Et ma scolarité en économie appliquée au C.T.P.E.A est complétée et attestée par la pièce jointe, mais le diplôme reste en attente de la soutenance de mon mémoire de sortie.

Je reste à votre entière disposition pour un entretien.

\makeletterclosing

\end{document}
